\section*{Note on Constraint Separation and the No-Refolding Condition}

We report a numerical result demonstrating the separation between two distinct constraints
on Hilbert-space dynamics: \emph{no-signaling} and a stronger constraint we call
\emph{no-refolding}. Central to this result is a precise operational definition of
``existing matter,'' formalized via a quadratic (fermionic) anchor.

\subsection*{Hilbert Space and Matter Anchor}

We consider a finite Hilbert space
\[
\mathcal{H} = (\mathbb{C}^2)^{\otimes N}, \quad N = 8,
\]
and an initial Hamiltonian $H_0$ describing fermionic matter.
Specifically, $H_0$ is a free-fermion hopping Hamiltonian defined on a one-dimensional
ring graph,
\[
H_0 = -t \sum_{\langle i,j\rangle \in A} \left( c_i^\dagger c_j + c_j^\dagger c_i \right),
\]
where $c_i^\dagger, c_i$ are Jordan--Wigner fermionic creation and annihilation operators
defined with respect to a fixed site ordering, and $A$ denotes the edge set of the ring.

This Hamiltonian is exactly quadratic in the fermionic operators and therefore defines
a stable fermionic matter sector. We take preservation of this quadratic structure as
our operational criterion for the persistence of matter.

\paragraph{Quadratic Anchor.}
For a general Hamiltonian $H$, we extract its quadratic component via a Hilbert--Schmidt
projection onto the fermionic bilinear operator basis $\{c_i^\dagger c_j\}$:
\[
t_{ij} \equiv
\frac{\mathrm{Tr}\!\left[(c_i^\dagger c_j)^\dagger H\right]}
     {\mathrm{Tr}\!\left[(c_i^\dagger c_j)^\dagger (c_i^\dagger c_j)\right]} .
\]
Using these coefficients, we reconstruct the best-fit quadratic Hamiltonian
\[
H_{\mathrm{quad}} = \sum_{i,j} t_{ij}\, c_i^\dagger c_j .
\]

We then define the \emph{quadratic fraction}
\[
Q(H) \equiv 1 - \frac{\|H - H_{\mathrm{quad}}\|}{\|H\|},
\]
where $\|\cdot\|$ denotes the Frobenius norm. By construction,
$Q(H)=1$ if and only if $H$ is exactly quadratic.

In our simulations, $Q(H_0)=1$ to numerical precision. We use $Q(H)$ as a
\emph{matter anchor}: values $Q(H)\approx 1$ indicate intact fermionic matter,
while $Q(H)\ll 1$ indicate destruction of quadratic (particle-like) structure.

\subsection*{Locality and Target Geometries}

Locality is defined relative to a chosen geometry via a signaling (locality) proxy.
Given a target graph $B$ on the same $N$ sites, we define a local operator subspace
consisting of on-site Pauli operators and two-body Pauli products supported on the edges
of $B$. The fraction of $\|H\|^2$ captured by this subspace defines the local weight,
and we denote the complementary nonlocal fraction as
\[
\mathrm{leak}_B(H) \equiv 1 - \frac{\|P_B(H)\|^2}{\|H\|^2}.
\]

As target geometries, we use connected random sparse graphs, which generically share
few or no edges with the original ring.

\subsection*{Dynamics}

We attempt to reduce $\mathrm{leak}_B(H)$ using only local unitary dynamics: conjugation
by adjacent two-qubit gates. This move set is identical in all runs and respects
no-signaling throughout.

We compare two dynamical regimes:
\begin{enumerate}
  \item \textbf{Free refolding}: minimize $\mathrm{leak}_B(H)$ without constraints.
  \item \textbf{No-refolding enforced}: minimize $\mathrm{leak}_B(H)$ subject to the
        hard constraint
        \[
        Q(H) \ge 0.98
        \]
        at every step of the evolution.
\end{enumerate}

The second regime enforces preservation of existing fermionic matter.

\subsection*{Result}

For certain target graphs, including cases where the target edge set is disjoint from
the source ring, we observe the following:

\begin{itemize}
  \item In the \emph{free} regime, $\mathrm{leak}_B(H)$ can be reduced substantially
        (e.g.\ from $1.0$ to $\sim 0.79$) using only local unitary gates.
        However, this reduction is accompanied by a catastrophic loss of quadratic
        structure, with final values $Q(H)\sim 10^{-3}$.
  \item In the \emph{no-refolding} regime, every proposed local update violates the
        quadratic constraint and is rejected. The dynamics stalls entirely:
        $\mathrm{leak}_B(H)$ remains equal to its initial value, while $Q(H)=1$ is
        preserved exactly.
\end{itemize}

Both regimes satisfy no-signaling identically; the sole difference is the requirement
that existing matter structure be preserved.

\subsection*{Constraint Separation}

These observations establish the following empirical fact:

\begin{quote}
There exist directions in Hilbert space that reduce signaling cost relative to a new
geometry but are dynamically accessible under local evolution \emph{only if} existing
fermionic (quadratic) structure is destroyed. When quadratic structure is preserved,
those geometries become dynamically inaccessible.
\end{quote}

We therefore identify \emph{no-refolding} as an independent constraint on Hilbert-space
dynamics, distinct from no-signaling. Once stable matter exists, the set of geometries
reachable by local unitary evolution is strictly restricted.

\subsection*{Implication}

Not all unitarily equivalent factorizations of Hilbert space are physically accessible
once matter is present. Geometry becomes history-dependent: the presence of matter
constrains which refactorizations of the Hilbert substrate can be dynamically realized.

This provides a concrete mechanism by which pre-geometric matter constrains emergent
spatial structure and motivates further investigation of dimensional selection as a
compatibility condition between matter and geometry.
